% Written By Enoch Yu
% Downloaded from archive.enochyu.com

\documentclass[12pt]{report}
\usepackage{graphicx}
\usepackage{amsmath,amsthm,amssymb}
\usepackage[font=small,labelfont=bf]{caption}
\usepackage{tikz}
\usepackage{float}
\usepackage[margin=1in]{geometry}
\usepackage{gensymb}
\usepackage{hyperref}
\hypersetup{
    colorlinks=true,
    linkcolor=blue,
    filecolor=magenta,      
    urlcolor=cyan,
    pdftitle={Ancient File 2},
    pdfpagemode=FullScreen,
}
\usepackage{fancyhdr}
\pagestyle{fancy}
\fancyhead[R]{Enoch Yu}
\newtheorem{theorem}{Theorem}[section]
\newtheorem{lemma}[theorem]{Lemma}
\newtheorem{proposition}{Proposition}
\newtheorem{corollary}{Corollary}[theorem]
\newcommand{\verteq}{\rotatebox{90}{$\;\;=\;\;$}}


\title{Ancient File 2}
\author{Enoch Yu}
\date{6 April 2025}

\begin{document}

\maketitle

\begin{center}
\small\em
Hello! My math group and I made bunch of fun problems
in the past, and I wrote solutions for some problems.
Apparently, I lost the problem sets and only have the
solutions :( Here's what I have!
\end{center}

\chapter*{Problem 3}
\section*{Key Word: Root of Unity, Euler's Formula, Vieta's Formula}
\subsection*{\textbf{Case I:} $\mathbf{x^{13}=1}$}
According to the Root of Unity, the values of $x$ may be rewritten. Moreover, according to Vieta's Theorem, because the sum of all $x$ is $-\frac{0}{1}$, the sum of all $x$ is 0. \\
\begin{tabular}{ccccccc}
    $x_1$ & $=$ & $e^{\frac{0}{13}\pi i}$ & $=$ & $\cos{\frac{0\pi}{13}} +i\sin{\frac{0\pi}{13}}$\\
    + & & & & \\
    $x_2$ & $=$ & $e^{\frac{2}{13}\pi i}$ & $=$ & $\cos{\frac{2\pi}{13}} +i\sin{\frac{2\pi}{13}}$ \\
    + & & & & \\
    $x_3$ & $=$ & $e^{\frac{4}{13}\pi i}$ & $=$ & $\cos{\frac{4\pi}{13}} +i\sin{\frac{4\pi}{13}}$ \\
    \vdots & & & & \vdots \\
    $x_{11}$ & $=$ & $e^{\frac{20}{13}\pi i}$ & $=$ & $\cos{\frac{20\pi}{13}} +i\sin{\frac{20\pi}{13}}$ \\
    + & & & & \\
    $x_{12}$ & $=$ & $e^{\frac{22}{13}\pi i}$ & $=$ & $\cos{\frac{22\pi}{13}} +i\sin{\frac{22\pi}{13}}$ \\
    + & & & & \\
    $x_{13}$ & $=$ & $e^{\frac{24}{13}\pi i}$ & $=$ & $\cos{\frac{24\pi}{13}} +i\sin{\frac{24\pi}{13}}$ \\
    $\verteq$ & & & & \\
    0 & & & & \\
\end{tabular}
\\\\
The properties of Sine and Cosine functions states that the following equations satisfy. \\
\setlength{\tabcolsep}{10pt}
\renewcommand{\arraystretch}{1.5}
\begin{tabular}{cccccccccccc}
    $\cos{\frac{2\pi}{13}}$ & = & $\cos{\frac{24\pi}{13}}$ & & & & & & $i\sin{\frac{2\pi}{13}}$ & = & $-i\sin{\frac{24\pi}{13}}$ \\
    $\cos{\frac{4\pi}{13}}$ & = & $\cos{\frac{22\pi}{13}}$ & & & & & & $i\sin{\frac{4\pi}{13}}$ & = & $-i\sin{\frac{22\pi}{13}}$ \\
    \vdots & = & \vdots & & & & & & \vdots & = & \vdots \\
    $\cos{\frac{10\pi}{13}}$ & = & $\cos{\frac{16\pi}{13}}$ & & & & & & $i\sin{\frac{10\pi}{13}}$ & = & $-i\sin{\frac{16\pi}{13}}$ \\
    $\cos{\frac{12\pi}{13}}$ & = & $\cos{\frac{14\pi}{13}}$ & & & & & & $i\sin{\frac{12\pi}{13}}$ & = & $-i\sin{\frac{14\pi}{13}}$
\end{tabular} \\
Therefore, the sum of solutions for the equation $x^{13}=1$ may be rewritten.
\begin{align*}
&x_1+x_2+\dots+x_{12}+x_{13}\\
=&
\cos{\frac{0\pi}{13}} +i\sin{\frac{0\pi}{13}}+
\cos{\frac{2\pi}{13}} +i\sin{\frac{2\pi}{13}}+
\dots +
\cos{\frac{22\pi}{13}} +i\sin{\frac{22\pi}{13}}+
\cos{\frac{24\pi}{13}} +i\sin{\frac{24\pi}{13}} \\
=&
\cos{\frac{0\pi}{13}} +i\sin{\frac{0\pi}{13}}+
\left( \cos{\frac{2\pi}{13}}+\dots+\cos{\frac{24\pi}{13}} \right)+
\left( i\sin{\frac{2\pi}{13}}+\dots+i\sin{\frac{24\pi}{13}} \right) \\
=&
\cos{0} +i\sin{0}+
2\left( \cos{\frac{14\pi}{13}}+\dots+\cos{\frac{24\pi}{13}} \right)+
(0) \\
=&
1 + 0 + 2\sum_{k=1}^{6}\cos\frac{(2k+12)\pi}{13} + 0 \\
=&
1 + 2\sum_{k=1}^{6}\cos\frac{(2k+12)\pi}{13} \\
\end{align*}
\begin{align*}
1 + 2\sum_{k=1}^{6}\cos\frac{(2k+12)\pi}{13}&=0 \\
\boldsymbol{\sum_{k=1}^{6}\cos\frac{(2k+12)\pi}{13}}&\boldsymbol{=-\frac{1}{2}}
\end{align*}

\subsection*{\textbf{Case II:} $\mathbf{x^{13}=-1}$}
According to Vieta's Theorem, because the sum of all $x$ is $-\frac{0}{1}$, the sum of all $x$ is 0. \\
\begin{tabular}{ccccccc}
    $x_1$ & $=$ & $e^{\frac{1}{13}\pi i}$ & $=$ & $\cos{\frac{1\pi}{13}} +i\sin{\frac{1\pi}{13}}$\\
    + & & & & \\
    $x_2$ & $=$ & $e^{\frac{3}{13}\pi i}$ & $=$ & $\cos{\frac{3\pi}{13}} +i\sin{\frac{3\pi}{13}}$ \\
    + & & & & \\
    $x_3$ & $=$ & $e^{\frac{5}{13}\pi i}$ & $=$ & $\cos{\frac{5\pi}{13}} +i\sin{\frac{5\pi}{13}}$ \\
    \vdots & & & & \vdots \\
    $x_3$ & $=$ & $e^{\frac{13}{13}\pi i}$ & $=$ & $\cos{\frac{13\pi}{13}} +i\sin{\frac{13\pi}{13}}$ \\
    \vdots & & & & \vdots \\
    $x_{11}$ & $=$ & $e^{\frac{21}{13}\pi i}$ & $=$ & $\cos{\frac{21\pi}{13}} +i\sin{\frac{21\pi}{13}}$ \\
    + & & & & \\
    $x_{12}$ & $=$ & $e^{\frac{23}{13}\pi i}$ & $=$ & $\cos{\frac{23\pi}{13}} +i\sin{\frac{23\pi}{13}}$ \\
    + & & & & \\
    $x_{13}$ & $=$ & $e^{\frac{25}{13}\pi i}$ & $=$ & $\cos{\frac{25\pi}{13}} +i\sin{\frac{25\pi}{13}}$ \\
    $\verteq$ & & & & \\
    0 & & & & \\
\end{tabular}
\\\\
The properties of Sine and Cosine functions states that the following equations satisfy. \\
\setlength{\tabcolsep}{10pt}
\renewcommand{\arraystretch}{1.5}
\begin{tabular}{cccccccccccc}
    $\cos{\frac{1\pi}{13}}$ & = & $\cos{\frac{25\pi}{13}}$ & & & & & & $i\sin{\frac{1\pi}{13}}$ & = & $-i\sin{\frac{25\pi}{13}}$ \\
    $\cos{\frac{3\pi}{13}}$ & = & $\cos{\frac{23\pi}{13}}$ & & & & & & $i\sin{\frac{3\pi}{13}}$ & = & $-i\sin{\frac{23\pi}{13}}$ \\
    \vdots & = & \vdots & & & & & & \vdots & = & \vdots \\
    $\cos{\frac{9\pi}{13}}$ & = & $\cos{\frac{17\pi}{13}}$ & & & & & & $i\sin{\frac{9\pi}{13}}$ & = & $-i\sin{\frac{17\pi}{13}}$ \\
    $\cos{\frac{11\pi}{13}}$ & = & $\cos{\frac{15\pi}{13}}$ & & & & & & $i\sin{\frac{11\pi}{13}}$ & = & $-i\sin{\frac{15\pi}{13}}$
\end{tabular} \\
Therefore, the sum of solutions for the equation $x^{13}=1$ may be rewritten.
\begin{align*}
&x_1+x_2+\dots+x_{12}+x_{13} \\
=&
\cos{\frac{1\pi}{13}} +i\sin{\frac{1\pi}{13}}+
\cos{\frac{3\pi}{13}} +i\sin{\frac{3\pi}{13}}+
\dots +
\cos{\frac{23\pi}{13}} +i\sin{\frac{23\pi}{13}}+
\cos{\frac{25\pi}{13}} +i\sin{\frac{25\pi}{13}} \\
=&
\cos{\frac{13\pi}{13}} +i\sin{\frac{13\pi}{13}}+
\left( \cos{\frac{1\pi}{13}}+\dots+\cos{\frac{11\pi}{13}}+\cos{\frac{15\pi}{13}}+\dots+\cos{\frac{25\pi}{13}} \right)+\\
&\text{ }\text{ }\text{ }\text{ }\text{ }\text{ }\text{ }\text{ }\text{ }\text{ }\text{ }\text{ }\text{ }\text{ }\text{ }\text{ }\text{ }\text{ }\text{ }\text{ }\text{ }\text{ }\text{ }\text{ }\text{ }\text{ }\text{ }\text{ }\text{ }\text{ }\text{ }\text{ }\text{ }\text{ }\text{ }\text{ }\text{ }\text{ }\text{ }\text{ }\text{ }\text{ }\text{ }\text{ }\text{ }\text{ }\text{ }\text{ }\text{ }\text{ }\text{ }\text{ }\text{ }\text{ }\text{ }\text{ }\left( i\sin{\frac{1\pi}{13}}+\dots+i\sin{\frac{11\pi}{13}}+i\sin{\frac{15\pi}{13}}+\dots+i\sin{\frac{25\pi}{13}} \right) \\
=&
\cos{\pi} +i\sin{\pi}+
2\left( \cos{\frac{1\pi}{13}}+\dots+\cos{\frac{11\pi}{13}} \right)+
(0) \\
=&
-1 + 0 + 2\sum_{k=1}^{6}\cos\frac{(2k-1)\pi}{13} + 0 \\
=&
-1 + 2\sum_{k=1}^{6}\cos\frac{(2k-1)\pi}{13} \\
\end{align*}
\begin{align*}
-1 + 2\sum_{k=1}^{6}\cos\frac{(2k-1)\pi}{13}&=0 \\
\boldsymbol{\sum_{k=1}^{6}\cos\frac{(2k-1)\pi}{13}}&\boldsymbol{=\frac{1}{2}}
\end{align*}
\subsection*{\textbf{Final Calculation}}
\[
\sum_{k=1}^{6}\cos\frac{(2k+12)\pi}{13}{=-\frac{1}{2}}
\]
\[
\sum_{k=1}^{6}\cos\frac{(2k-1)\pi}{13}=\frac{1}{2}
\]
\[
\sum_{k=1}^{6}\cos\frac{(2k+12)\pi}{13}\cdot\sum_{k=1}^{6}\cos\frac{(2k-1)\pi}{13}=-\frac{1}{2}\cdot\frac{1}{2}=-\frac{1}{4}
\]
\\\\
\\\\
Conclusion: \textbf{$-\frac{1}{4}$ is the calculated value}.


\chapter*{Problem 4}
\section*{Key Word: Pigeon Hole Principle}
Let $D_n$ be the accumulated number of completed homework problems on $n^{th}$ day.
\\\\
Because students must solve at least one problem each day and 100 in total, the following expression is feasible.
\begin{equation}
1\le D_1 < D_2 < \dots < D_{64} < D_{65} = 100
\end{equation}
A positive integer $x$ can be arbitrarily added in the expression to provide the following:
\begin{equation}
1+x\le D_1 +x < D_2 +x < \dots < D_{64}+x < D_{65}+x = 100+x.
\end{equation}
The expressions (1) and (2) may be combined.
\begin{equation}
1\le D_1, D_2, \dots, D_{64}, D_{65}, D_1 +x, D_2 +x, \dots, D_{64}+x < D_{65}+x = 100+x
\end{equation}
According to expression (1), it is evident that $D_1, D_2, \dots, D_{64}, D_{65}$ are all different. Similarly, $D_1 +x, D_2 +x, \dots, D_{64}+x, D_{65}+x$ are all distinct according to expressions (2).
\\\\
There exists $130$ pigeons for expression (3) and retains $100+x$ pigeon holes. According to the Pigeon Hole Principle, if the number of pigeons exceeds the number of pigeon holes, at least two pigeons must share a hole. In another term, if $130>100+x$, then at least two numbers from $D_1, D_2, \dots, D_{64}, D_{65}, D_1 +x, D_2 +x, \dots, D_{64}+x, D_{65}+x$ must be same because pigeon hole represents quantity in this instance. However, if two values are to be same, $D_a=D_b+x$ because $D_a\ne D_b$ and $D_a+x \ne D_b+x$.
\\\\
If $130>100+x$, then $D_a=D_b+x$, or there always exists two days in which total of $x$ problems are solved during those two days. Because $x<30$, the maximum possible value for $x$ is $29$.
\\\\
\\\\
Conclusion: \textbf{29 is the maximum possible value of $x$}.



\chapter*{Problem 5}
\section*{Key Word: Lucas's Theorem, Stars and Bars, Property of Modular}
\begin{tabular}{c|c|c|c}
    \boxed{1246} & \boxed{896} & \boxed{4050} & \boxed{7547} \\
    \text{1246 Stars, 57 Bars} & \text{896 Stars, 57 Bars} & \text{4050 Stars, 57 Bars} & \text{7547 Stars, 57 Bars} \\
    $\frac{(1246+57)!}{1246!57!}$ & $\frac{(896+57)!}{896!57!}$ & $\frac{(4050+57)!}{4050!57!}$ & $\frac{(7547+57)!}{7547!57!}$ \\
    $={}_{1303}C_{57}$ & $={}_{953}C_{57}$ & $={}_{4107}C_{57}$ & $={}_{7604}C_{57}$ \\
\end{tabular}
\\\\
The total probability is ${}_{1303}C_{57} \cdot {}_{953}C_{57} \cdot {}_{4107}C_{57} \cdot {}_{7604}C_{57}$

\section*{Applying Lucas's Theorem for $r_7$}
\begin{align*}
57_{(10)}&=111_{(7)} \\
1303_{(10)}&=3542_{(7)} \\
953_{(10)}&=2531_{(7)} \\
4107_{(10)}&=14655_{(7)} \\
7604_{(10)}&=31111_{(7)} \\
\end{align*}
Lucas's Theorem is applied to find $r_7$.
\begin{align*}
    {}_{1303}C_{57}
    &\equiv {}_{3}C_{0}\cdot{}_{5}C_{1}\cdot{}_{4}C_{1}\cdot{}_{2}C_{1} \pmod{7} \\
    &\equiv 1\cdot5\cdot4\cdot2 \pmod{7} \\
    &\equiv 40 \pmod{7} \\
    &\equiv 5 \pmod{7}
    \\\\
    {}_{953}C_{57}
    &\equiv {}_{2}C_{0}\cdot{}_{5}C_{1}\cdot{}_{3}C_{1}\cdot{}_{1}C_{1} \pmod{7} \\
    &\equiv 1\cdot5\cdot3\cdot1 \pmod{7} \\
    &\equiv 15 \pmod{7} \\
    &\equiv 1 \pmod{7}
    \\\\
    {}_{4107}C_{57}
    &\equiv {}_{1}C_{0}\cdot{}_{4}C_{0}\cdot{}_{6}C_{1}\cdot{}_{5}C_{1}\cdot{}_{5}C_{1} \pmod{7} \\
    &\equiv 1\cdot1\cdot6\cdot5\cdot5 \pmod{7} \\
    &\equiv 150 \pmod{7} \\
    &\equiv 3 \pmod{7}
    \\\\
    {}_{7604}C_{57}
    &\equiv {}_{3}C_{0}\cdot{}_{1}C_{0}\cdot{}_{1}C_{1}\cdot{}_{1}C_{1}\cdot{}_{1}C_{1} \pmod{7} \\
    &\equiv 1\cdot1\cdot1\cdot1\cdot1 \pmod{7} \\
    &\equiv 1 \pmod{7}  
\end{align*}
$r_7 \Rightarrow 5\cdot1\cdot3\cdot1=15\Rightarrow1 \text{ }(\because15 \equiv1 \pmod{7})$
$\therefore r_7=1$

\section*{Applying Lucas's Theorem for $r_3$}
\begin{align*}
57_{(10)}&=2010_{(3)} \\
1303_{(10)}&=1201021_{(3)} \\
953_{(10)}&=1022022_{(3)} \\
4107_{(10)}&=111212100_{(3)} \\
7604_{(10)}&=101102122_{(3)} \\
\end{align*}
Lucas's Theorem is applied to find $r_7$.
\begin{align*}
    {}_{1303}C_{57}
    &\equiv {}_{1}C_{0}\cdot{}_{2}C_{0}\cdot{}_{0}C_{0}\cdot{}_{1}C_{2}\cdot{}_{0}C_{0}\cdot{}_{2}C_{1}\cdot{}_{1}C_{0} \pmod{3} \\
    &\equiv 1\cdot1\cdot1\cdot0\cdot1\cdot2\cdot1 \pmod{3} \\
    &\equiv 0 \pmod{3}
\end{align*}
$r_3=0 \text{ }(\because {}_{1303}C_{57} \equiv0 \pmod{3})$ \\
$\therefore r_3=0$
\\\\
$r_7 \cdot r_3=1\cdot0=0$
\\\\
\\\\
Conclusion: \textbf{$r_7 \cdot r_3$ is equal to zero}.



\chapter*{Problem 10}
\section*{Key Word: Cauchy-Schwarz Inequality}
\section*{Part I}
According to Cauchy-Schwarz inequality, for any number $a_n$ and $b_n$, \\$({a_0}^2+\dots+{a_n}^2)({b_0}^2+\dots+{b_n}^2)\ge(a_0b_0+\dots+a_nb_n)^2$ is true.
\\\\
$\frac{4}{3}x^2+5y^2+\frac{1}{2}z^2=3$ and $f(x,y,z)=3x+2y-6z$ are given by the problem. In order to find the maximum value for $f(x,y,z)$, Cauchy-Schwarz inequality can be utilized.
\[
\left\{
\left(\sqrt{\frac{4}{3}x^2}\right)^2
+\left(\sqrt{5y^2}\right)^2
+\left(\sqrt{\frac{1}{2}z^2}\right)^2
\right\}
\left\{
\left(\frac{3}{\sqrt{\frac{4}{3}}}\right)^2
+\left(\frac{2}{\sqrt{5}}\right)^2
+\left(\frac{-6}{\sqrt{\frac{1}{2}}}\right)^2
\right\}
\ge(3x+2y-6z)^2
\]
The Cauchy-Schwarz inequality was written to utilize $\frac{4}{3}x^2+5y^2+\frac{1}{2}z^2=3$ and the fact that the maximum value of $3x+2y-6z$ must be found. The inequality may be simplified.
\[
\left\{
\left(\frac{4}{3}x^2\right)
+\left(5y^2\right)
+\left(\frac{1}{2}z^2\right)
\right\}
\left\{
\left(\frac{9}{\frac{4}{3}}\right)
+\left(\frac{4}{5}\right)
+\left(\frac{36}{\frac{1}{2}}\right)
\right\}
\ge(3x+2y-6z)^2
\]
Substituting $\frac{4}{3}x^2+5y^2+\frac{1}{2}z^2$ provides the maximum value of $(3x+2y-6z)^2$.
\[
\{3\}
\left\{
\left(\frac{27}{4}\right)+\left(\frac{4}{5}\right)+(72)
\right\}
\ge(3x+2y-6z)^2
\]
\[
\{3\}
\left\{
\left(\frac{135}{20}\right)+\left(\frac{16}{20}\right)+\left(\frac{1440}{20}\right)
\right\}
\ge(3x+2y-6z)^2
\]
\[
\{3\}
\left\{\frac{1591}{20}\right\}
\ge(3x+2y-6z)^2
\]
\[
\frac{4773}{20}
\ge(3x+2y-6z)^2
\]
By square rooting both sides, the maximum value of $3x+2y-6z$ could be found.
\[
-\sqrt{\frac{4773}{20}}\le\sqrt{(3x+2y-6z)^2}\le\sqrt{\frac{4773}{20}}
\]
\[
-\sqrt{\frac{4773}{20}}\le3x+2y-6z\le\sqrt{\frac{4773}{20}}
\]
\[
-\frac{1}{2}\sqrt{\frac{4773}{5}}\le3x+2y-6z\le\frac{1}{2}\sqrt{\frac{4773}{5}}
\]
\\\\
\\\\
Conclusion: \textbf{The maximum value of $f(x,y,z)$ is $\frac{1}{2}\sqrt{\frac{4773}{5}}$}.


\section*{Part II}
\begin{figure}[H]
    \includegraphics[scale=0.15]{Visualization.png}
    \caption{GeoGebra was utilized to manifest the current situation.}
\end{figure}
$3x+2y-6z=\sqrt{\frac{4773}{20}}$ is the blue plane that was established in Part I. By investigating the points where the plane meets $x-axis, y-axis \text{ and } z-axis$, the area formula for tetrahedron is available to discover the volume of the polyhedron enclosed by this plane and the constraints $x \ge 0, y \ge 0, \text{ and } z \le 0$.

\subsection*{Intersection with $x$-axis}
$y=0, z=0$ \\
$3x+2(0)-6(0)=\sqrt{\frac{4773}{20}}$ \\
$3x=\sqrt{\frac{4773}{20}}$ \\
$\therefore x=\frac{1}{3}\sqrt{\frac{4773}{20}}$
\\\\
$(\frac{1}{3}\sqrt{\frac{4773}{20}},0,0)$ is the intersection point.

\subsection*{Intersection with $y$-axis}
$x=0, z=0$ \\
$3(0)+2y-6(0)=\sqrt{\frac{4773}{20}}$ \\
$2y=\sqrt{\frac{4773}{20}}$ \\
$\therefore y=\frac{1}{2}\sqrt{\frac{4773}{20}}$
\\\\
$(0,\frac{1}{2}\sqrt{\frac{4773}{20}},0)$ is the intersection point.

\subsection*{Intersection with $z$-axis}
$x=0, y=0$ \\
$3(0)+2(0)-6z=\sqrt{\frac{4773}{20}}$ \\
$-6z=\sqrt{\frac{4773}{20}}$ \\
$\therefore z=-\frac{1}{6}\sqrt{\frac{4773}{20}}$
\\\\
$(0,0,-\frac{1}{6}\sqrt{\frac{4773}{20}})$ is the intersection point.

\subsection*{Finding the Volume}
\begin{align*}
    \text{Volume} &= \frac{1}{3}\cdot\left|\frac{1}{3}\sqrt{\frac{4773}{20}}\right|\cdot\left|\frac{1}{2}\sqrt{\frac{4773}{20}}\right|\cdot\left|-\frac{1}{6}\sqrt{\frac{4773}{20}}\right| \\
    &= \frac{1}{3}\cdot\frac{1}{6}\cdot\frac{4773}{20}\cdot\frac{1}{6}\cdot\sqrt{\frac{4773}{5}} \\
    &=\frac{1}{3}\cdot\frac{1}{2}\cdot\frac{1591}{20}\cdot\frac{1}{6}\cdot\sqrt{\frac{4773}{5}} \\
    &=\frac{1591}{720}\cdot\frac{1}{2}\sqrt{\frac{4773}{5}} \\
    &=\frac{1591}{1440}\sqrt{\frac{4773}{5}} \\
    &=\frac{1591\sqrt{23865}}{7200}
\end{align*}

\vspace{\baselineskip}
\vspace{\baselineskip}
\noindent
Conclusion: \textbf{The volume of the polyhedron enclosed by the plane and the constraints $x \ge 0$, $y \ge 0$ and $z \le 0$ is $\frac{1591\sqrt{23865}}{7200}$}.
\\\\
\\\\
\noindent
Conclusion: \textbf{The maximum value for $f(x,y,z)$ is $\frac{\sqrt{23865}}{10}$. The volume of the polyhedron enclosed by this plane and the constraints $x \ge 0$, $y \ge 0$ and $z \le 0$ is $\frac{1591\sqrt{23865}}{7200}$}.



\chapter*{Problem 11}
\section*{Key Word: Vieta's Theoream, AM-GM Inequality}
The generalized form for AM-GM inequality states that $a_1+a_2+\dots+a_n \ge n\sqrt[n]{a_1a_2\dots a_n}$ for $a_n>0$.
\\\\
According to Vieta's theorem, the following equations are true.
\begin{align*}
    -a &= r_1+r_2+\dots+r_8 \\
    -b &= r_1r_2r_3++\dots+r_6r_7r_8 \\
    -c &= r_1r_2r_3r_4r_5++\dots+r_4r_5r_6r_7r_8 \\
    -d &= r_1r_2r_3r_4r_5r_6r_7++\dots+r_2r_3r_4r_5r_6r_7r_8 \\
    e &= r_1r_2r_3r_4r_5r_6r_7r_8 \\
\end{align*}
The equations above are symmetric. Meaning, the following expression is true using AM-GM inequality.
\begin{multline*}
-(a+b+c+d)=(r_1+\dots+r_8)+(r_1r_2r_3+\dots+r_6r_7r_8)+(r_1\dots r_5+\dots+r_4\dots r_8)+\\(r_1\dots r_7+\dots+r_2\dots r_8) \ge n\sqrt[n]{r_1\dots r_8)^m}
\end{multline*}
$n$, which represents the number of terms in the left side of the inequality, could be investigated using combination.
\begin{align*}
    n &={_8}C_1+{_8}C_3+{_8}C_5+{_8}C_7 \\
    &=8+56+56+8 \\
    &= 128
\end{align*}
Using the symmetric nature of $-a,-b,-c\text{ and }-d$, it is evident that the number of $r_1$ equals the number of $r_1,r_2,r_3,r_4,r_5,r_6,r_7 \text{ and }r_8$. In another terms, $m$ is equal to the number of $r_1$. $m$ could be found by assuming that $r_1$ is already chosen and choosing either 0, 2, 4 or 6 terms from $r_2,r_3,r_4,r_5,r_6 \text{ and } r_7$ because 1, 3, 5 or 7 terms are selected in total from $r_1,r_2,r_3,r_4,r_5,r_6,r_7 \text{ and }r_8$.
\begin{align*}
    m &={_7}C_0+{_7}C_2+{_7}C_4+{_7}C_6 \\
    &=1+21+35+7 \\
    &= 64
\end{align*}
The following inequality may be written by substituting the known values.
\[
-(a+b+c+d)\ge 128\sqrt[128]{e^{64}}
\]
Because both $-(a+b+c+d)e$ and $128\sqrt[128]{e^{64}}$ are positive numbers, the entire inequality continues to satisfy when both sides are squared.
\[
\left(-(a+b+c+d)\right)^2\ge \left(128\sqrt[128]{e^{64}}\right)^2
\]
\[
(a+b+c+d)^2\ge \left(128e^{\frac{1}{2}}\right)^2
\]
\[
(a+b+c+d)^2\ge 2^{14}e
\]
\[
\frac{(a+b+c+d)^2}{e} \ge 2^{14}
\]
The least possible value of $\frac{(a+b+c+d)^2}{e}$ is $2^{14}$, or 16384.
\\\\
\\\\
Conclusion: \textbf{The minimum possible value for $\frac{(a+b+c+d)^2}{e}$ is 16384}.


\chapter*{Problem 14}
\section*{Key Word: Mass Point}
\includegraphics[]{Screenshot 2025-04-06 at 8.57.09 PM.png}
\end{document}


